% META: {"id": "1NSI-TC-EX-023", "chapitre": "1NSI-TYPES-CONSTRUITS", "type_objet": "exercice", "capacites": ["1NSI-TYPES-CONSTRUITS-C4"], "parcours": 1, "competences": ["programmer"], "duree_min": 8, "mode_creation": "ex_nihilo", "sources_inspiration": [], "status": "approved", "capacites_codes": ["C4"], "methodes": ["M1"], "fichier_tex": "chapitres/1NSI-TYPES-CONSTRUITS/exercices/1NSI-TC-EX-023.tex", "corrige_tex": "chapitres/1NSI-TYPES-CONSTRUITS/corriges/1NSI-TC-CO-023.tex"}
% BEGIN-VERIFY
% def fusionner(tab1, tab2):
%     resultat = []
%     i = 0
%     j = 0
%     while i < len(tab1) and j < len(tab2):
%         if tab1[i] <= tab2[j]:
%             resultat.append(tab1[i])
%             i = i + 1
%         else:
%             resultat.append(tab2[j])
%             j = j + 1
%     while i < len(tab1):
%         resultat.append(tab1[i])
%         i = i + 1
%     while j < len(tab2):
%         resultat.append(tab2[j])
%         j = j + 1
%     return resultat
% assert fusionner([1, 3, 5], [2, 4, 6]) == [1, 2, 3, 4, 5, 6]
% assert fusionner([], [1, 2]) == [1, 2]
% assert fusionner([1], []) == [1]
% assert fusionner([1, 1], [1, 1]) == [1, 1, 1, 1]
% END-VERIFY
\begin{exercice}{1NSI-TC-EX-023}{3}{15}
On dispose de deux tableaux d'entiers, chacun trié par ordre croissant.

Écrire une fonction \lstinline{fusionner(tab1, tab2)} qui renvoie un nouveau tableau
contenant tous les éléments des deux tableaux, trié par ordre croissant, sans
utiliser de fonction de tri.

\textbf{Exemple :}
\begin{console}
>>> fusionner([1, 3, 5], [2, 4, 6])
[1, 2, 3, 4, 5, 6]
\end{console}

\textit{Indication :} utiliser deux indices parcourant chacun des tableaux.
\end{exercice}
